\providecommand{\ZZ}{\mathbb{Z}}
\providecommand{\CC}{\mathbb{C}}
\providecommand{\cO}{\mathcal{O}}
\providecommand{\Kum}{\mathrm{Kum}}
\providecommand{\DT}{\mathrm{DT}}
\providecommand{\Eone}{E_1}
\providecommand{\Etwo}{E_2}
\providecommand{\kBKM}{\kappa_{\mathrm{BKM}}}
\providecommand{\kch}{\kappa_{\mathrm{ch}}}
\providecommand{\kcat}{\kappa_{\mathrm{cat}}}
\providecommand{\kfib}{\kappa_{\mathrm{fiber}}}
\providecommand{\PhiN}{\Phi_N}

\section*{The order-$8$ Nikulin boundary row}

Let $g_8$ be a symplectic automorphism of order $8$ of a K3 surface.
In the order-indexed Nikulin-Gritsenko normalisation, Borcherds'
singular-theta weight formula applied to the twined weak Jacobi input
gives
\[
  c_8(0)=2,\qquad \kBKM(\Phi_8)=\frac{c_8(0)}{2}=1.
\]
This is an automorphic weight statement. A compact Calabi-Yau threefold
source, a Donaldson-Thomas theory, and a chiral Hall algebra require
independent geometric data.

\subsection*{Three row systems}

The primitive compact CHL ladder is
\[
  N \in \{1,2,3,4,6\},
  \qquad
  (c_N(0),\kBKM(\Phi_N))
  =
  (10,5),(8,4),(6,3),(4,2),(2,1).
\]
These five orders are exactly the orders compatible with the
origin-preserving automorphism group of an elliptic curve in the compact
diagonal quotient of $K3\times E$.

The order-indexed Nikulin-Mukai extension is
\[
\begin{array}{c|ccc}
N & 5 & 7 & 8\\ \hline
c_N(0) & 4 & 2 & 2\\
\kBKM(\Phi_N)=c_N(0)/2 & 2 & 1 & 1 .
\end{array}
\]
The entries $N=5,7,8$ are automorphic boundary rows. A compact
CY$_3$ host, a reduced $\DT$ theory, and a CY-to-chiral realisation are
additional data.

The Gritsenko-Clery simplest-divisor catalogue is indexed by triples
$(t,N;k)$:
\[
\begin{gathered}
 (1,1;5),\ (2,1;2),\ (1,2;3),\ (3,1;1),\\
 (1,3;2),\ (4,1;\tfrac12),\ (1,4;\tfrac32),\ (2,2;1).
\end{gathered}
\]
Its weight tuple is
\[
  (5,2,3,1,2,\tfrac12,\tfrac32,1)
\]
in the displayed order. Its common entry with the primitive CHL ladder
is $\Delta_5$. Using the scalar weight $1$ of the triple
$(2,2;1)$ as order-$8$ CY$_3$ data requires an explicit row map
specifying the input Jacobi form, multiplier, paramodular group, and
constant-term normalisation. The symbol \(N=2\) in the triple
$(2,2;1)$ belongs to the diagonal-divisor indexing, separate from the
primitive CHL order and from the order-$8$ Nikulin row.

\subsection*{The level-one product}

The compact product $K3\times E$ supplies the level-one denominator
data:
\[
  \kcat(K3\times E)=0,\qquad
  \kch(K3\times E)=0,\qquad
  \kch^{\mathrm{Heis}}(K3\times E)=3,
\]
\[
  \kBKM(\mathfrak g_{\Delta_5})=5,\qquad
  \kfib(K3)=24.
\]
The normalised denominator is
\[
  D_5=64^{-1}\Delta_5,\qquad \Phi_{10}=\Delta_5^2 .
\]
These numbers answer different questions. The numerical expression
$5=3+2$ combines three different constructions: \(3\) is the
Heisenberg-Mukai specialisation, \(2=\chi(\cO_{K3})\) is the K3 fibre
Euler characteristic, and \(5\) is the Borcherds weight of \(\Delta_5\).
The compact $K3\times E$ Hodge supertrace is \(0\), and the CHL ladder
reads the Borcherds coefficient \(c_N(0)/2\).

At order $8$, the primitive CHL construction is outside compact diagonal
quotients of $K3\times E$: the elliptic factor has origin-preserving
automorphism groups of orders \(2\), \(4\), or \(6\), hence element
orders in \(\{1,2,3,4,6\}\). A translation quotient or stacky quotient is
a separate geometric package and must carry its own reduced $\DT$
orientation, partition identity, and Borcherds root-multiplicity
theorem.

\subsection*{Generalised Kummer dimension}

For a complex two-torus $A$, Beauville's generalised Kummer variety is
the fibre
\[
  \Kum_n(A)
  =
  \sigma^{-1}(0_A)
  \subset A^{[n+1]},
  \qquad
  \sigma(\xi)=\sum_{x\in \xi}x ,
\]
and has complex dimension $2n$. Hence the generalised Kummer fourfold
is $\Kum_2(A)$, the Albanese fibre of $A^{[3]}\to A$. The fibre with
$n=3$ has complex dimension $6$. Any notation of the form
$\Kum^3(A)$ must therefore specify its convention before it can be used
as fourfold data.

Mongardi-Tari-Wandel, \emph{Automorphisms of generalised Kummer
fourfolds}, arXiv:1512.00225, classifies finite symplectic automorphism
groups of Kummer fourfolds by lattice methods. In Beauville's
convention this is the fourfold \(\Kum_2(A)\); \(\Kum_3(A)\) is the
sixfold. The result is group and lattice boundary data for hyperkahler
fourfolds of Kummer type. A compact CY$_3$ source for $\Phi_8$ requires
a separate construction.

\subsection*{Kummer quotients and Calabi-Yau host data}

Let $X=\Kum_2(A)$ be a fourfold of Kummer type and let
\(h\in\operatorname{Aut}(X)\) be a symplectic automorphism of order
\(8\). The quotient stack \([X/\langle h\rangle]\) has dimension \(4\)
and \(\omega\cong\cO\) in the stacky sense. Its coarse quotient
\(X/\langle h\rangle\) has symplectic quotient singularities. If a
crepant symplectic resolution exists, the resolved space remains
holomorphic symplectic and carries a nonzero holomorphic two-form. This
is irreducible holomorphic symplectic geometry, while a strict
Calabi-Yau fourfold has \(h^{0,2}=0\).

The Bridgeland-King-Reid derived McKay theorem supplies equivalences
under its own hypotheses. A local quotient
\[
  \CC^4/\ZZ_8(1,3,5,7)
\]
is Gorenstein because $1+3+5+7\equiv 0\pmod 8$, and the weights pair as
\((1,7)\) and \((3,5)\) for the symplectic form. This is a local
canonical-class and symplectic linear calculation. Smoothness of the
$G$-Hilbert scheme, existence of a global crepant resolution,
orientation data for a reduced enumerative theory, and compatibility
with the Borcherds product remain separate theorems.

A compact or stacky host \(X_8\) requires fresh computation of
$\kcat$ and $\kch$. For irreducible holomorphic symplectic manifolds of
dimension $2n$ one has
$h^{0,2i}=1$ for $0\leq i\leq n$, hence
\[
  \chi(\cO)=n+1
\]
on the smooth hyperkahler model. Thus \(\Kum_2(A)\) has
\(\chi(\cO)=3\). A strict smooth Calabi-Yau threefold has compact Hodge
supertrace \(1-1=0\), and a strict smooth Calabi-Yau fourfold has
compact Hodge supertrace \(1+1=2\). These Hodge numbers are geometric
invariants of candidate hosts, distinct from the $N=8$ Borcherds weight
\(1\).

\subsection*{The positive half}

The automorphic side supplies the weight
\[
  \kBKM(\Phi_8)=1.
\]
The BPS Lie bracket, Hall product, and chiral algebra require their own
construction. For a compact or stacky host $X_8$, the positive half must
be constructed as
\[
  Y^+(X_8)
  =
  H^\bullet_{\mathrm{eq}}
  \bigl(\mathcal M^+_{\mathrm{eff}}(X_8),\phi_W\bigr)
\]
with the exact equivariance, orientation, and obstruction theory
specified. Since $d\geq 3$, the specialised chiral output is
\Eone-chiral. The \Etwo-braided object belongs to the Drinfeld centre
or to a completed Hall-Drinfeld double after the required pairing and
coproduct are constructed.

The Hall-Drinfeld branch attached to the Borcherds product reads root
multiplicities from the Fourier coefficients of the chosen input. The
Mukai Yangian branch is the self-mirror K3 object on the Mukai lattice.
The CoHA positive half is $Y^+$; in the basic local test case
\(\operatorname{CoHA}(\CC^3)=Y^+\). The full
$\mathcal W_{1+\infty}$ comparison belongs to a Drinfeld double,
derived centre, completion, or vertex-level evaluation setting.

\subsection*{Automorphic theorem and geometric construction}

The proved statement is automorphic:
\[
  \boxed{\;\kBKM(\Phi_8)=c_8(0)/2=1\;}
\]
for the order-$8$ Nikulin boundary row. The source data are the
symplectic K3 automorphism of order $8$, the frame shape
$1^2 2^1 4^1 8^2$, and the Gritsenko-Nikulin correction giving
$c_8(0)=2$. This row lies outside the primitive compact CHL ladder.
Identifying it with the Gritsenko-Clery triple \((2,2;1)\) requires a
separate row map.

A geometric realisation requires:
\begin{itemize}
\item a compact or stacky Calabi-Yau threefold $X_8$ carrying the
      order-$8$ datum;
\item a reduced $\DT$ obstruction theory and orientation compatible
      with that datum;
\item a weak Jacobi input identical to the one used for $c_8(0)=2$;
\item a partition identity whose automorphic side is $\Phi_8$ in the
      order-indexed normalisation;
\item a Hall or BPS bracket theorem matching the Fourier coefficients
      of the same Borcherds product;
\item for any Kummer-type candidate, an explicit declaration that the
      object is a hyperkahler fourfold, a stacky Calabi-Yau object, a
      local quotient model, or a compact CY$_3$, together with its
      quotient action, singularity type, crepant-resolution existence, and
      Hodge calculation.
\end{itemize}
Mongardi-Tari-Wandel supplies boundary evidence for finite symplectic
group actions in Kummer-fourfold geometry. This evidence controls the
order range; the CY$_3$ package above remains additional data.

\subsection*{References}

Beauville 1983 defines the generalised Kummer varieties and their
dimension. Nikulin 1980 and Mukai 1988 classify the finite symplectic
K3 orders. Borcherds 1998 gives the singular-theta weight formula.
Gritsenko-Nikulin and Gritsenko-Clery give the paramodular
Borcherds-weight inputs in their separate indexing systems.
Mongardi-Tari-Wandel 2015 classifies automorphisms of generalised
Kummer fourfolds.
